\begin{center}
	\section*{Homework 11 - 3.12}
	Due May 9 \\
	Uzair Hamed Mohammed
\end{center}

% 4, 8, 14, 16, 20, 22, 24

\subsection*{3.12 - Moment-Generating Functions}
\begin{enumerate}
	\item[4.] Find the moment-generating function for the discrete random variable X whose probability function is given by
	      \[
		      p_X (k) = (\frac{3}{4})^k (\frac{1}{4}), k = 0, 1, 2, \dots
	      \]

	      \underline{Sol}: \\

	      The probability function is that of a geometric distribution with success probability \(p = \frac{1}{4}\) counting failures:
	      \[
		      p_X(k) = \left(\frac{3}{4}\right)^k \frac{1}{4}, \quad k = 0, 1, 2, \dots
	      \]
	      The moment-generating function is
	      \[
		      M_X(t) = \operatorname{E}[e^{tX}] = \sum_{k=0}^\infty e^{tk} \left(\frac{3}{4}\right)^k \frac{1}{4}
		      = \frac{1}{4} \sum_{k=0}^\infty \left(\frac{3e^t}{4}\right)^k.
	      \]
	      The series converges for \(\frac{3e^t}{4} < 1\), i.e., \(t < \ln\frac{4}{3}\). Summing gives
	      \[
		      M_X(t) = \frac{1}{4} \cdot \frac{1}{1 - \frac{3}{4}e^t} = \frac{1}{4 - 3e^t},
		      \qquad t < \ln\frac{4}{3}.
	      \]

	\item[8.] Let Y be a continuous random variable with \(f_Y (y) = ye^{-y}, 0 \leq y\). Show that \(M_Y (t) = \frac{1}{(1 - t)^2}\).

	      \underline{Sol}: \\

	      \[
		      M_Y(t) = \operatorname{E}[e^{tY}] = \int_0^\infty e^{ty} y e^{-y} \, dy = \int_0^\infty y e^{-(1-t)y} \, dy, \quad t<1.
	      \]
	      Let \(u = (1-t)y\), \(du = (1-t)dy\), \(y = u/(1-t)\):
	      \[
		      M_Y(t) = \int_0^\infty \frac{u}{1-t} e^{-u} \frac{du}{1-t} = \frac{1}{(1-t)^2} \int_0^\infty u e^{-u} \, du = \frac{\Gamma(2)}{(1-t)^2} = \frac{1}{(1-t)^2}.
	      \]

	\item[14.] Find an expression for \(E(Y^k)\) if \(M_Y (t) = (1 - t / \lambda)^{- \lambda}\), where \(\lambda\) is any positive real number and \(r\) is a positive integer.

	      \underline{Sol}: \\

	      \[
		      M_Y(t) = \left(1 - \frac{t}{\lambda}\right)^{-\lambda} \quad\Rightarrow\quad Y \sim \text{Gamma}(\alpha=\lambda,\ \beta=\lambda).
	      \]
	      \[
		      \operatorname{E}[Y^k] = \frac{\lambda^\lambda}{\Gamma(\lambda)} \int_0^\infty y^{\lambda+k-1} e^{-\lambda y}\,dy
		      = \frac{\lambda^\lambda}{\Gamma(\lambda)} \cdot \frac{\Gamma(\lambda+k)}{\lambda^{\lambda+k}}
		      = \frac{\Gamma(\lambda+k)}{\Gamma(\lambda)\,\lambda^k}.
	      \]
	      Thus
	      \[
		      \boxed{\operatorname{E}[Y^k] = \frac{(\lambda)_k}{\lambda^k} = \prod_{i=0}^{k-1}\left(1+\frac{i}{\lambda}\right)}.
	      \]

	\item[16.] Find the variance Y is \(M_Y (t) = e^{2t} / (1 - t^2)\).

	      \underline{Sol}: \\

	      \[
		      M_Y(t) = \frac{e^{2t}}{1-t^2},\quad |t|<1.
	      \]

	      \[
		      M_Y(t) = \left(1 + 2t + \frac{(2t)^2}{2!} + \frac{(2t)^3}{3!} + \cdots\right)\left(1 + t^2 + t^4 + \cdots\right).
	      \]

	      Coefficient of \(t\): \(2\) \(\Rightarrow\) \(E[Y]=2\).

	      Coefficient of \(t^2\): from \( \frac{(2t)^2}{2!}=2t^2 \) and \(1\cdot t^2 = t^2 \) gives \(3t^2\) \(\Rightarrow\) \(\frac{E[Y^2]}{2!}=3\) \(\Rightarrow\) \(E[Y^2]=6\).

	      \[
		      \operatorname{Var}(Y)=6-2^2=2.
	      \]

	\item[20.] Calculate \(P(X \leq 2)\) if \(M_X (t) = (\frac{1}{4} + \frac{3}{4}e^t)^5\).

	      \underline{Sol}: \\

	      \[
		      M_X(t) = \left(\frac14 + \frac34 e^t\right)^5 \implies X \sim \text{Binomial}(n=5, p=\tfrac34).
	      \]
	      \[
		      P(X \le 2) = \sum_{k=0}^2 \binom{5}{k} \left(\frac34\right)^k \left(\frac14\right)^{5-k}
		      = \frac{1}{4^5} \sum_{k=0}^2 \binom{5}{k} 3^k.
	      \]
	      \[
		      \binom{5}{0}3^0 = 1,\quad \binom{5}{1}3^1 = 15,\quad \binom{5}{2}3^2 = 10\cdot 9 = 90.
	      \]
	      Sum = \(1+15+90 = 106\). \(4^5 = 1024\). \[
		      P(X \le 2) = \frac{106}{1024} = \frac{53}{512}.
	      \]

	\item[22.] Suppose the moment-generating function for a random variable W is given by
	      \[
		      M_W (t) = e^{-3 + 3e^t} \cdot (\frac{2}{3} + \frac{1}{3} e^t)^4
	      \]

	      Calculate \(P(W \leq 1).\) (Hint: Write W as a sum.)

	      \underline{Sol}: \\

	      \[
		      M_W(t)=e^{-3+3e^t}\left(\frac{2}{3}+\frac{1}{3}e^t\right)^4
	      \]
	      implies \(W = X + Y\) with \(X\sim\text{Poisson}(3)\), \(Y\sim\text{Binomial}(4,\frac13)\), independent.

	      \[
		      P(W=0)=P(X=0)P(Y=0)=e^{-3}\left(\frac{2}{3}\right)^4 = e^{-3}\cdot\frac{16}{81}.
	      \]
	      \[
		      \begin{aligned}
			      P(W=1)=P(X=1)P(Y=0)+P(X=0)P(Y=1) \\
			      =3e^{-3}\cdot\frac{16}{81}+e^{-3}\cdot\binom{4}{1}\frac13\left(\frac23\right)^3
			      = \frac{48}{81}e^{-3} + e^{-3}\cdot\frac{32}{81}= \frac{80}{81}e^{-3}.
		      \end{aligned}
	      \]
	      \[
		      P(W\le 1)=\left(\frac{16}{81}+\frac{80}{81}\right)e^{-3}= \frac{96}{81}e^{-3}= \frac{32}{27}e^{-3}.
	      \]

	\item[24.] Suppose that Y is a normal variable, where \(f_Y (y) = (1 / \sqrt{2 \pi} \sigma) \textrm{ exp } \lbrack - \frac{1}{2} (\frac{y - \mu}{\sigma})^2 \rbrack, - \infty < y < \infty\).
	      \begin{enumerate}
		      \item Does the random variable \(W = 3Y\) have a normal distribution?
		      \item DOes the random variable \(W = 3Y + 1\) have a normal distribution?
	      \end{enumerate}

	      \underline{Sol}: \\

	      \[
		      Y \sim N(\mu,\sigma^2) \implies aY+b \sim N(a\mu+b, a^2\sigma^2).
	      \]
	      \begin{enumerate}
		      \item \(W=3Y\): \(a=3,b=0\) \(\Rightarrow\) normal.
		      \item \(W=3Y+1\): \(a=3,b=1\) \(\Rightarrow\) normal.
	      \end{enumerate}
\end{enumerate}
